\begin{helper}[Right-side bridge response]
Let $t,u$ be the right-side data of a bridge, and let $x$ be a current
left-side set with $\mathsf{InitialSegment}(x,u)$.  Under either legal right
response from \noderef{FrontBridgeStep}, the updated data satisfy
the following common conclusions: $t'$ belongs to the prefix tree and is
nonempty, $u'$ is nonempty, and for the response point $q$ we have
\[
  j<q\quad(j\in u),\qquad u'=\operatorname{insert}(q,u).
\]
Moreover $\mathsf{InitialSegment}(x,u')$; if $t'\in F$ then
$\mathsf{InitialSegment}(t',\mathsf{FrontBridgeTail}(u'))$, while if
$t'\notin F$ then $t'=\mathsf{FrontBridgeTail}(u')$.
\end{helper}
\begin{proof}
In the nonterminal right response, write the new point as $n$.  The
finite-tail insertion lemma \noderef{FrontBridgeTailInsert}, together with
$t=\mathsf{FrontBridgeTail}(u)$, gives
\[
 \mathsf{FrontBridgeTail}(u')
 =\operatorname{insert}(n,\mathsf{FrontBridgeTail}(u))
 =\operatorname{insert}(n,t)=t'.
\]
Thus the two alternatives for the final tail condition follow respectively
from equality and from the same equality.  The given tree membership and
nonemptiness of $t'$ are retained, and $u'$ is nonempty because it contains
the old nonempty set $u$.  Applying \noderef{InitialSegmentInsert} to
$\mathsf{InitialSegment}(x,u)$ gives the updated left initial segment.

In the terminal right response, write the new point as $k$.  The same tail
insertion lemma applies.  Since every element of $\mathsf{FrontBridgeTail}(u)$
is an element of $u$, every such element is below $k$.  Hence
\noderef{InitialSegmentInsert}, applied to
$\mathsf{InitialSegment}(t,\mathsf{FrontBridgeTail}(u))$, yields
\[
 \mathsf{InitialSegment}(t',\mathsf{FrontBridgeTail}(u')).
\]
Here $t'=t$, so its tree membership and nonemptiness come from the old
bridge.  The updated left initial segment follows by applying the same
lemma to $\mathsf{InitialSegment}(x,u)$.  Finally, the response hypothesis
ensures $t'=t\in F$, so the alternative $t'\notin F$ is impossible.  In
both cases the displayed response point supplies the asserted existential
witness for $q$.
\end{proof}
